\section{Conclusion}
\label{sec:conclusion}

We asked how much small, meaning-preserving edits to an instruction change LLM task performance. To answer it, we built a controlled protocol: six benign edit types, a seeded editor whose combined edits reuse the local edits of their parts, and a paired item-level analysis with flip rates and additive and union baselines. Across two small instruction-tuned models, four tasks and 20,800 scored generations, a single edit typically moves scores by 0--5 points, but specific edit--task pairs move them by 10--25 points. The damaging edit is not a fixed edit type. Edits that alter the answer-space specification, such as the punctuation and casing of a label list, break classification. Edits that weaken the output-format constraint break strict-format multiple choice and, when combined, generation. Generation was not the most fragile task type, and super-additive compounding appeared only there. Elsewhere, effects saturate because an item can break only once.

{\bf Key takeaway:} prompt fragility follows what an edit touches, not what kind of edit it is. Label lists and output-format instructions are load-bearing parts of a prompt. Sensitivity metrics should be reported together with the parser and token budget used, and with flip rates rather than only net changes.

\para{Future work.} The most important next step is to repeat the study with frontier models, which only requires pointing our code at a hosted endpoint. We also plan to use several base templates per task, to separate edit type from template structure, and to run more seeds per edit type. Adding a lenient or LLM-judge scorer to every condition would separate format and content fragility throughout. Longer-form generation, such as summarization and chain-of-thought math, would test whether compounding grows with output length. Finally, a targeted experiment that edits only the label list or only the task description would directly test the answer-space mechanism.
